% Original Prism starter using the official Springer Nature class.
% Initial draft; see Prism docs/guide.zh-CN.md for scope and final-file requirements.
% The official class and bibliography style are imported during initialization.
\documentclass[pdflatex,sn-nature]{sn-jnl}
\usepackage{graphicx,amsmath,amssymb,amsfonts,booktabs,array,tabularx,siunitx}
\sisetup{uncertainty-mode=separate}

\begin{document}
\title[Short manuscript title]{Manuscript title}
\author*[1]{\fnm{First} \sur{Author}}\email{first.author@example.org}
\author[1,2]{\fnm{Second} \sur{Author}}
\affil*[1]{\orgdiv{Department},\orgname{First institution},
  \orgaddress{\city{City},\country{Country}}}
\affil[2]{\orgdiv{Department},\orgname{Second institution},
  \orgaddress{\city{City},\country{Country}}}
\abstract{Replace this paragraph with an unreferenced abstract describing
the research problem, approach, verified findings and their implications.
For a Nature Machine Intelligence Article, the current guidance allows up
to 150 words. All examples in this starter demonstrate typesetting only;
they are not scientific results.}
\maketitle

% NMI Articles begin with introductory text without an Introduction heading.
Replace this introductory text with the context, unresolved question and
specific contribution of the study. Explain the problem for readers outside
the immediate field. The citation below tests bibliography mechanics only;
retain it only if it is relevant to the manuscript\cite{wilkinson2016fair}.

This is a working manuscript, using the official Nature reference style.
It does not reproduce the published journal's final layout. Replace every
instructional statement, synthetic number and graphic slot before submission.

\section{Results}\label{sec:results}
\subsection{Primary comparison}
State the verified primary result and connect it to the evidence needed to
assess it. Figure~\ref{fig:comparison} illustrates related panels under one
legend; Table~\ref{tab:comparison} illustrates uncertainty alignment.

\begin{figure}[htbp]
\centering
\begin{minipage}[t]{0.47\linewidth}
\raggedright\textbf{a}\par\smallskip
\fbox{\parbox[c][35mm][c]{0.90\linewidth}{\centering
Overview graphic\\[5pt]Replace with a PDF or PNG}}
\end{minipage}\hfill
\begin{minipage}[t]{0.47\linewidth}
\raggedright\textbf{b}\par\smallskip
\fbox{\parbox[c][35mm][c]{0.90\linewidth}{\centering
Comparison graphic\\[5pt]Use matching label sizes}}
\end{minipage}
\caption{\textbf{Two related panels in a single figure.} \textbf{a}, Describe
the overview. \textbf{b}, Describe the comparison. Define the actual sample
sizes, units, symbols and error bars from the study. The boxes are synthetic
layout slots; no scientific conclusion is represented.}
\label{fig:comparison}
\end{figure}

\begin{table}[htbp]
\centering
\caption{Synthetic numerical examples for table layout only.}
\label{tab:comparison}
\small
\begin{tabularx}{\linewidth}{@{}>{\raggedright\arraybackslash}X
  S[table-format=3.0] S[table-format=1.3(3)]@{}}
\toprule
Setting & {Count} & {Estimate $\pm$ SD}\\
\midrule
Example A & 120 & 0.385 +- 0.006\\
Example B & 98 & 0.608 +- 0.051\\
\bottomrule
\end{tabularx}
\par\smallskip\footnotesize\raggedright
All values are invented. State the actual meaning of the estimate and its
uncertainty when replacing the cells. Do not infer a scientific claim from
this formatting example.
\end{table}

\subsection{Additional analysis}
Present secondary analyses and validation in the order needed to assess the
main finding. Use topical subheadings. Keep main display items focused and
move detailed supporting material to the independent supplementary document.

\section{Discussion}\label{sec:discussion}
Interpret the verified findings, their relation to prior work, limitations
and implications. For NMI Articles, the Discussion does not use subheadings.
Distinguish evidence from speculation and avoid repeating the Results.

\section{Methods}\label{sec:methods}
\subsection{Study design}
Describe materials, data sources, inclusion rules, controls and versions.
Include the relevant ethical or reporting details supplied by the authors.

\subsection{Mathematical formulation}
Define each variable and its units before use. Equations~\eqref{eq:objective}
and~\eqref{eq:activation} illustrate grouped notation; replace them with the
study's justified equations or remove them.
\begin{subequations}\label{eq:model}
\begin{align}
  \mathcal{L}(\theta)&=\frac{1}{n}\sum_{i=1}^{n}
    \ell\bigl(f_\theta(x_i),y_i\bigr)+\lambda\lVert\theta\rVert_2^2,
    \label{eq:objective}\\
  h(z)&=\begin{cases}z,&z\geq0,\\\alpha z,&z<0.\end{cases}
    \label{eq:activation}
\end{align}
\end{subequations}
These are notation examples, not assumptions or conclusions about a real
study. Refer to a numbered expression using \verb|\eqref{eq:objective}|.

\subsection{Validation and statistical analysis}
Describe the unit of analysis, exact sample counts, independent repetitions,
uncertainty measures and validation procedures. Report statistical tests and
their assumptions only where they were actually used.

\section*{Data availability}
Replace with an accurate statement of where the study's data can be accessed,
including stable identifiers and any access restrictions.

\section*{Code availability}
Replace with an accurate statement of code availability, versions and access
conditions. Do not publish private code or credentials through this template.

\section*{Acknowledgements}
Replace with confirmed acknowledgements, or remove if not applicable.

\section*{Author contributions}
Replace with the actual contributions agreed by all authors.

\section*{Funding}
Replace with confirmed funding sources and grant identifiers, or the accurate
statement appropriate to the work.

\section*{Competing interests}
Replace with the authors' confirmed competing-interest statement.

% The sn-nature class option already selects \bibliographystyle{sn-nature}.
\bibliography{references}
\end{document}
